\paragrafos{Διαίρεση κλασμάτων: πόσες φορές χωράει}{Α.2.6, «Διαίρεση κλασμάτων»}

\begin{anagke}
Ο παππούς Γιάννης ετοιμάζει δείγματα ελαιόλαδου για το παζάρι. Έχει $\frac{3}{4}$ του λίτρου λάδι και θέλει να το βάλει σε μικρά μπουκαλάκια που χωράνε το καθένα $\frac{1}{8}$ του λίτρου. «Πόσα μπουκαλάκια θα γεμίσουν;» ρωτάει, και η Ζωή και ο Ηλίας αρχίζουν να σκέφτονται.

Ο Ηλίας γράφει τα $\frac{3}{4}$ ως $\frac{6}{8}$. «Το κάθε μπουκαλάκι είναι ένα όγδοο. Έχω έξι όγδοα. Έξι μπουκαλάκια». Η Ζωή ρωτάει: «Αυτό είναι διαίρεση. Πώς γράφεται;» Γράφουν $\frac{3}{4}:\frac{1}{8}$. Ο Στέλιος, που θέλει πάντα έναν κανόνα, προσπαθεί: «Διαιρώ αριθμητή με αριθμητή και παρονομαστή με παρονομαστή. Το $3:1=3$ και το $4:8$... δεν βγαίνει φυσικός». Η Ρήνη σηκώνει τους ώμους: «Βγάζεις περίεργα νούμερα».

Η Πολυξένη παρατηρεί κάτι στην απάντηση του Ηλία: έξι είναι ακριβώς $\frac{3}{4}\cdot 8$. «Μήπως η διαίρεση με το $\frac{1}{8}$ είναι πολλαπλασιασμός με το 8, δηλαδή με τον αντίστροφό του;» Αν είναι έτσι, ο κανόνας υπάρχει και δεν χρειάζεται καν τα όγδοα.
\end{anagke}

\begin{lexi}{σύνθετο κλάσμα}{σύνθετο κλάσμα}
Η λέξη \lexiword{σύνθετο} σημαίνει «φτιαγμένο από πολλά κομμάτια». Σύνθετο κλάσμα είναι ένα κλάσμα που έχει μέσα του άλλα κλάσματα, στον αριθμητή ή στον παρονομαστή ή και στα δύο, για παράδειγμα $\dfrac{\frac{2}{3}}{\frac{4}{9}}$. Η κλασματική γραμμή του σημαίνει διαίρεση: $\frac{2}{3}:\frac{4}{9}$.

Και η διαίρεση, στη γλώσσα της ζωής, απαντάει στο \emph{πόσες φορές χωράει}: πόσες φορές χωράει το μικρό μπουκαλάκι στο λάδι που έχω.
\end{lexi}

\begin{metafrasi}
Πόσες φορές χωράει το $\frac{1}{8}$ στα $\frac{3}{4}$; & $\dfrac{3}{4}:\dfrac{1}{8}=\dfrac{3}{4}\cdot\dfrac{8}{1}=6$\\
Διαίρεση σημαίνει πολλαπλασιασμός με τον αντίστροφο. & $\dfrac{\alpha}{\beta}:\dfrac{\gamma}{\delta}=\dfrac{\alpha}{\beta}\cdot\dfrac{\delta}{\gamma}$\\
Ο φυσικός 6 διαιρούμενος με το $\frac{2}{3}$. & $6:\dfrac{2}{3}=6\cdot\dfrac{3}{2}=9$\\
Το μισό από τα $\frac{3}{4}$. & $\dfrac{3}{4}:2=\dfrac{3}{4}\cdot\dfrac{1}{2}=\dfrac{3}{8}$\\
Σύνθετο κλάσμα. & $\dfrac{\frac{2}{3}}{\frac{5}{7}}=\dfrac{2}{3}:\dfrac{5}{7}=\dfrac{14}{15}$\\
\end{metafrasi}

\tmima{Η θεωρία}

Στους φυσικούς, για να διαιρέσουμε το 6 με το 3, ρωτάμε πόσες φορές χωράει το 3 στο 6. Το ίδιο ισχύει και στα κλάσματα. Και υπάρχει ένας πολύ βολικός κανόνας.

\begin{kanonas}{διαίρεση}
\begin{itemize}[leftmargin=1.2em,itemsep=0.3ex,topsep=2pt]
\item Για να διαιρέσω δύο φυσικούς: πολλαπλασιάζω τον διαιρετέο με τον αντίστροφο του διαιρέτη. $\alpha:\beta=\alpha\cdot\dfrac{1}{\beta}$.
\item Για να διαιρέσω δύο κλάσματα: πολλαπλασιάζω τον διαιρετέο με τον \textbf{αντίστροφο του διαιρέτη}. $\dfrac{\alpha}{\beta}:\dfrac{\gamma}{\delta}=\dfrac{\alpha}{\beta}\cdot\dfrac{\delta}{\gamma}$.
\end{itemize}
\end{kanonas}

Γιατί; Στο λάδι του παππού: $\frac{3}{4}:\frac{1}{8}$ ρωτάει πόσες φορές χωράει το $\frac{1}{8}$ στα $\frac{3}{4}$. Σε ένα ολόκληρο λίτρο χωράνε 8 τέτοια μπουκαλάκια (γιατί $8\cdot\frac{1}{8}=1$). Στα $\frac{3}{4}$ του λίτρου χωράνε τα $\frac{3}{4}$ από το 8, δηλαδή $\frac{3}{4}\cdot 8=6$. Ο αριθμός 8 είναι ο αντίστροφος του $\frac{1}{8}$.

\tmimasmall{Σύνθετο κλάσμα}

Ένα κλάσμα που έχει κλάσμα στον αριθμητή ή στον παρονομαστή λέγεται \textbf{σύνθετο}. Η κλασματική του γραμμή σημαίνει διαίρεση, άρα το μετατρέπω σε απλό με τον κανόνα της διαίρεσης:
$$\dfrac{\frac{\alpha}{\beta}}{\frac{\gamma}{\delta}}=\dfrac{\alpha}{\beta}:\dfrac{\gamma}{\delta}=\dfrac{\alpha\cdot\delta}{\beta\cdot\gamma}.$$
Αν ο διαιρετέος ή ο διαιρέτης είναι μεικτός αριθμός, τον γράφω πρώτα ως κλάσμα.

\begin{prosoxi}
\begin{itemize}[leftmargin=1.2em,itemsep=0.3ex,topsep=2pt]
\item Αντιστρέφουμε \textbf{μόνο τον διαιρέτη}, όχι τον διαιρετέο. $\frac{3}{4}:\frac{1}{8}=\frac{3}{4}\cdot 8=6$, ενώ ο «αντίστροφος και των δύο», $\frac{4}{3}\cdot 8=\frac{32}{3}$, είναι λάθος.
\item Η διαίρεση με κλάσμα μικρότερο από το 1 \textbf{μεγαλώνει} τον αριθμό: $6:\frac{1}{2}=12$, γιατί στο 6 χωράνε 12 μισά. Η διαίρεση με φυσικό μεγαλύτερο από το 1 τον μικραίνει.
\item Δεν διαιρούμε με το 0, όπως πάντα. Ούτε με το κλάσμα $\frac{0}{5}$, που είναι 0.
\end{itemize}
\end{prosoxi}

\tmima{Λυμένα παραδείγματα}

\paradeigma{Τα δείγματα του παππού}
\tixero Ο παππούς Γιάννης έχει $\frac{3}{4}$ του λίτρου λάδι και το βάζει σε μπουκαλάκια των $\frac{1}{8}$ του λίτρου.\par
\tithelo Πόσα μπουκαλάκια γεμίζουν.\par
\possindeo Ζητάω πόσες φορές χωράει το $\frac{1}{8}$ στο $\frac{3}{4}$, δηλαδή $\frac{3}{4}:\frac{1}{8}$. Ο αντίστροφος του $\frac{1}{8}$ είναι το 8. Άρα $\frac{3}{4}\cdot 8=\frac{24}{4}=6$.\par
\apantisi 6 μπουκαλάκια.\par
\elegxos $6\cdot\frac{1}{8}=\frac{6}{8}=\frac{3}{4}$. Βγαίνει το λάδι που είχαμε.

\paradeigma{Μπουκαλάκια που δεν χωράνε ακριβώς}
\tixero Ο παππούς Γιάννης έχει $\frac{2}{3}$ του λίτρου και μπουκαλάκια των $\frac{1}{4}$ του λίτρου.\par
\tithelo Πόσα μπουκαλάκια θα γεμίσουν.\par
\possindeo $\frac{2}{3}:\frac{1}{4}=\frac{2}{3}\cdot 4=\frac{8}{3}=2\frac{2}{3}$.\par
\apantisi Γεμίζουν 2 ολόκληρα μπουκαλάκια και τα $\frac{2}{3}$ ενός τρίτου.\par
\elegxos $2\cdot\frac{1}{4}=\frac{1}{2}$ και μένουν $\frac{2}{3}-\frac{1}{2}=\frac{1}{6}$ του λίτρου. Το $\frac{1}{6}$ είναι τα $\frac{1}{6}:\frac{1}{4}=\frac{4}{6}=\frac{2}{3}$ του μπουκαλιού. Συμφωνεί.

\paradeigma{Σύνθετα κλάσματα}
\tixero Τα κλάσματα $\dfrac{\frac{5}{6}}{\frac{10}{9}}$ και $\dfrac{4}{\frac{2}{5}}$.\par
\tithelo Να τα κάνω απλά.\par
\possindeo Η κλασματική γραμμή σημαίνει διαίρεση. $\dfrac{\frac{5}{6}}{\frac{10}{9}}=\dfrac{5}{6}:\dfrac{10}{9}=\dfrac{5}{6}\cdot\dfrac{9}{10}=\dfrac{45}{60}=\dfrac{3}{4}$. Και $\dfrac{4}{\frac{2}{5}}=4:\dfrac{2}{5}=4\cdot\dfrac{5}{2}=10$.\par
\apantisi $\dfrac{3}{4}$ και $10$.\par
\elegxos Το $\frac{3}{4}\cdot\frac{10}{9}=\frac{30}{36}=\frac{5}{6}$ και το $10\cdot\frac{2}{5}=4$. Σωστά.

\paradeigma{Διαίρεση μεικτών}
\tixero Οι αριθμοί $3\frac{1}{3}$ και $1\frac{1}{9}$.\par
\tithelo Το πηλίκο $3\frac{1}{3}:1\frac{1}{9}$.\par
\possindeo $3\frac{1}{3}=\frac{10}{3}$ και $1\frac{1}{9}=\frac{10}{9}$. Άρα $\frac{10}{3}:\frac{10}{9}=\frac{10}{3}\cdot\frac{9}{10}=\frac{\cancel{10}\cdot 9}{3\cdot\cancel{10}}=3$.\par
\apantisi 3.\par
\elegxos $3\cdot 1\frac{1}{9}=3\cdot\frac{10}{9}=\frac{10}{3}=3\frac{1}{3}$. Σωστά.

\tmima{Ασκήσεις}

\zoni{Το κατάλαβα}{εφαρμογή του κανόνα}

\ex Κάνε τις διαιρέσεις: α) $\frac{3}{5}:\frac{1}{10}$ \quad β) $\frac{5}{6}:\frac{5}{12}$ \quad γ) $\frac{7}{8}:\frac{3}{4}$ \quad δ) $\frac{4}{9}:\frac{8}{27}$ \quad ε) $\frac{2}{3}:2$ \quad στ) $6:\frac{3}{4}$.

\ex Κάνε τις διαιρέσεις: α) $12:\frac{2}{3}$ \quad β) $5:\frac{1}{4}$ \quad γ) $\frac{3}{7}:3$ \quad δ) $10:\frac{5}{2}$ \quad ε) $\frac{1}{2}:\frac{1}{6}$ \quad στ) $\frac{9}{10}:\frac{3}{5}$.

\ex Πόσες φορές χωράει: α) το $\frac{1}{4}$ στο 1 \quad β) το $\frac{1}{3}$ στα $\frac{2}{3}$ \quad γ) το $\frac{1}{6}$ στα $\frac{5}{12}$ \quad δ) το $\frac{2}{5}$ στο 2.

\ex Υπολόγισε με μεικτούς: α) $2\frac{1}{2}:\frac{5}{6}$ \quad β) $3\frac{1}{3}:2$ \quad γ) $4:1\frac{1}{3}$ \quad δ) $1\frac{3}{4}:2\frac{1}{2}$.

\ex Μετάτρεψε σε απλά τα σύνθετα κλάσματα: α) $\dfrac{\frac{2}{3}}{\frac{4}{9}}$ \quad β) $\dfrac{5}{\frac{5}{6}}$ \quad γ) $\dfrac{\frac{7}{8}}{7}$ \quad δ) $\dfrac{\frac{3}{4}}{\frac{9}{10}}$.

\ex Ποιο κλάσμα πολλαπλασιασμένο με το $\frac{3}{5}$ δίνει $\frac{9}{10}$;

\ex Βρες τον αριθμό, αν: α) το $\frac{1}{3}$ του είναι $\frac{7}{9}$ \quad β) τα $\frac{2}{5}$ του είναι $\frac{6}{7}$.

\ex Υπολόγισε με προτεραιότητα πράξεων: α) $\frac{1}{2}:\frac{1}{4}\cdot\frac{1}{3}$ \quad β) $\frac{1}{2}:\left(\frac{1}{4}\cdot\frac{1}{3}\right)$ \quad γ) $\left(\frac{3}{4}+\frac{1}{4}\right):\frac{1}{2}$ \quad δ) $\frac{3}{4}+\frac{1}{4}:\frac{1}{2}$.

\ex Σωστό ή λάθος; α) $\frac{3}{4}:\frac{1}{2}=\frac{3}{8}$ \quad β) $6:\frac{1}{3}=18$ \quad γ) Η διαίρεση με κλάσμα μικρότερο από το 1 μεγαλώνει τον αριθμό. \quad δ) $\frac{2}{5}:\frac{2}{5}=1$ \quad ε) $\frac{1}{2}:3=6$ \quad στ) Για να διαιρέσω αντιστρέφω και τα δύο κλάσματα.

\zoni{Το χρησιμοποιώ}{προβλήματα από τη ζωή}

\ex \textbf{Το λάδι.} Ο παππούς Γιάννης έχει $\frac{5}{6}$ του λίτρου λάδι και μπουκαλάκια των $\frac{1}{12}$ του λίτρου. Πόσα γεμίζει; Αν είχε μπουκαλάκια των $\frac{1}{5}$ του λίτρου, πόσα θα γέμιζε και πόσο μέρος του τελευταίου θα έμενε γεμάτο;

\ex \textbf{Η κορδέλα.} Η Αγγελική έχει μια κορδέλα μήκους $\frac{9}{2}$ μ. και την κόβει σε κομμάτια μήκους $\frac{3}{4}$ μ. Πόσα κομμάτια βγάζει;

\ex \textbf{Η ταχύτητα.} Ο Νίκος τρέχει $2\frac{1}{2}$ χλμ. σε $\frac{1}{4}$ της ώρας (15 λεπτά). Πόσα χιλιόμετρα θα έτρεχε σε μια ώρα με την ίδια ταχύτητα;

\ex \textbf{Οι βάρδιες.} Η θεία Μαρία έχει $3\frac{1}{2}$ ώρες δουλειά στην αποθήκη και κάθε βάρδια διαρκεί $\frac{7}{4}$ της ώρας. Πόσες βάρδιες χρειάζεται;

\ex \textbf{Οι μερίδες.} Η γιαγιά Μυρτώ έχει 3 πίτες και κόβει κάθε μερίδα στα $\frac{3}{8}$ μιας πίτας. Πόσες μερίδες βγάζει;

\ex \textbf{Η συνταγή.} Για 6 άτομα η συνταγή θέλει $\frac{3}{4}$ του κιλού αλεύρι. Πόσο αλεύρι αντιστοιχεί σε ένα άτομο;

\zoni{Το σκέφτομαι}{μια ερώτηση που πάει πιο μακριά}

\ex Ο Στέλιος γράφει $\frac{3}{4}:\frac{1}{8}=\frac{4}{3}\cdot 8=\frac{32}{3}$. Βρες το λάθος του και το σωστό αποτέλεσμα.

\ex Τι βγαίνει αν διαιρέσεις ένα κλάσμα με τον αντίστροφό του; Δοκίμασε με τα $\frac{2}{3}$ και με τα $\frac{3}{5}$.

\ex Γιατί το $6:\frac{1}{2}=12$ είναι μεγαλύτερο από το 6; Εξήγησε με το «πόσα μισά χωράνε στο 6».

\ex Οι Αιγύπτιοι έγραφαν τα $\frac{2}{3}$ του $\frac{1}{7}$ ως $\frac{1}{14}+\frac{1}{42}$. Έλεγξε ότι αυτό ισούται με $\frac{2}{3}\cdot\frac{1}{7}$. Βρες με τον ίδιο τρόπο το $\frac{2}{3}$ του $\frac{1}{5}$ ως άθροισμα δύο κλασματικών μονάδων.

\ex Ποιος αριθμός πολλαπλασιασμένος επί $\frac{4}{5}$ κάνει 12; Ποιος αριθμός διαιρεμένος με το $\frac{4}{5}$ κάνει 12;

\zoni{Για το σπίτι}{με το μυαλό και με το μπουκάλι}

\ex Πόσα ποτήρια των $\frac{1}{4}$ του λίτρου γεμίζουν από ένα μπουκάλι $1\frac{1}{2}$ λίτρου; Έλεγξε με πραγματικό μπουκάλι και ποτήρι.

\ex Υπολόγισε το σύνθετο κλάσμα $\dfrac{\frac{1}{2}+\frac{1}{3}}{\frac{1}{2}-\frac{1}{3}}$.
